% chapters/06_hierarchical_rl_architecture.tex

\section{Overview}

The convergence theory of Chapter 5 established that a carefully designed
three-timescale actor-critic algorithm can provably learn the mean-field game
equilibrium with common noise, provided the function approximators satisfy strong
regularity conditions. However, the practical implementation of such an algorithm in
high-dimensional financial markets requires a structured architecture that respects
the underlying mathematical decomposition of the problem. This chapter develops
exactly such an architecture.

The central insight --- formalised in Proposition~\ref{prop:timescale-separation} ---
is that the master equation of the MFG naturally decomposes the learning dynamics
into three distinct operators corresponding to distinct physical and economic
phenomena: idiosyncratic diffusion, mean-field interaction, and common-noise drift.
These operators operate on fundamentally different timescales and involve different
information structures. We show that this decomposition canonically motivates a
hierarchical reinforcement learning architecture comprising three specialised
components:

\begin{itemize}
    \item A \textbf{fast critic} (Soft Actor-Critic, SAC) that learns the value of
    individual states under idiosyncratic noise;
    \item An \textbf{intermediate actor} (Proximal Policy Optimization, PPO) that
    adjusts the policy in response to changes in the population distribution;
    \item A \textbf{slow meta-learner} (Twin Delayed Deep Deterministic Policy
    Gradient, TD3) that adapts to shifts in the common-noise regime.
\end{itemize}

\textbf{We emphasize that the specific assignment of SAC, PPO, and TD3 to the three
operators is a heuristic design choice guided by the SPDE structure, not a proven
optimal mapping.} The hierarchical principle --- fast critic, intermediate actor,
slow meta-learner --- is the fundamental contribution; the particular algorithms are
illustrative and can be replaced by any critic, actor, and meta-learner that operate
on distinct timescales and optimise objectives aligned with the SPDE decomposition.

The chapter provides a formal timescale separation principle
(Proposition~\ref{prop:timescale-separation}) that justifies the decoupling of these
three learning processes, and demonstrates --- via ablation studies on the
linear-quadratic benchmark and a simulated multi-asset market --- that the
hierarchical architecture significantly outperforms monolithic RL approaches in terms
of sample efficiency and convergence stability. \textbf{We include a control
experiment with randomly chosen timescale ratios to isolate the contribution of the
SPDE-guided timescale separation.}

\paragraph{Prior work.} The SPDE formulation of the master equation for common-noise
MFGs originates from \citep{carmonadelaruelachapelle2014}. The use of
multi-timescale RL for MFGs has been explored by \citep{angiulietal2024} for finite spaces, and by \citep{zhouzhouhu2025} for continuous-time
actor-critic flows. Zeng et al.\ (2024) proposed a single-loop algorithm with
finite-time guarantees. The present chapter contributes a concrete hierarchical
architecture (SAC-PPO-TD3) grounded in the SPDE operator decomposition and validated
on continuous-space benchmarks.

\paragraph{Chapter organisation.} \Cref{sec:hrl-decomposition} recalls the SPDE
decomposition of the master equation. \Cref{sec:hrl-mapping} maps each operator to a
specific algorithmic component, with explicit acknowledgment of the heuristic nature
of the mapping. \Cref{sec:hrl-proposition} states the formal timescale separation
principle (Proposition~\ref{prop:timescale-separation}) and discusses its
assumptions. \Cref{sec:hrl-implementation} describes the implementation of the
hierarchical architecture. \Cref{sec:hrl-experiments} presents experimental results,
including an ablation study that isolates the SPDE decomposition's contribution.
\Cref{sec:hrl-limitations} discusses limitations. \Cref{sec:hrl-conclusion}
concludes.


\section{SPDE Decomposition of the Master Equation}
\label{sec:hrl-decomposition}

Recall from Chapter 1 that the MFG equilibrium is characterised by the coupled FBSDE
system \eqref{eq:fbsde-system}. The value function $U(t,\mu)$ of the representative
agent, viewed as a function on the Wasserstein space, satisfies the master equation
(\citep{carmonadelaruelachapelle2014}
\[
    \partial_t U + \frac{\sigma^2}{2} \int_{\mathbb{R}^d} \Delta_x \frac{\delta
    U}{\delta\mu}(\mu)(x)\, \mu(dx) + \int_{\mathbb{R}^d} H^*(t,x,\mu,D_x U)\,
    \mu(dx) = 0,
\]
where $\frac{\delta U}{\delta\mu}$ is the Lions derivative (\Cref{sec:mp-wasserstein}).
In the presence of common noise, the measure flow $\mu_t$ is stochastic and evolves
according to the Fokker--Planck SPDE \eqref{eq:fp-spde-strong}. The master equation
can be decomposed into three distinct operators, each governing a different aspect of
the dynamics.

\paragraph{Operator 1: Idiosyncratic Diffusion (Generator $\mathcal{L}_{\mathrm{idio}}$).}
This operator captures the effect of the agent's own idiosyncratic noise on the value
function. It is given by
\begin{equation}
    \mathcal{L}_{\mathrm{idio}} U(t,\mu) = \frac{\sigma^2}{2} \int_{\mathbb{R}^d}
    \Delta_x \frac{\delta U}{\delta\mu}(\mu)(x)\, \mu(dx).
    \label{eq:operator-idio}
\end{equation}
This term involves only the agent's private state and the idiosyncratic volatility
$\sigma$. It operates on the fastest timescale because idiosyncratic shocks are
frequent and independent across agents. Learning this operator corresponds to
estimating the value of individual states and the local curvature of the value
function.

\paragraph{Operator 2: Mean-Field Interaction (Drift Derivative $\mathcal{L}_{\mathrm{mf}}$).}
The mean-field interaction enters through the Hamiltonian $H^*$, specifically through
the dependence of the drift $b$ and the reward $f$ on the population measure $\mu$.
The corresponding operator is
\begin{equation}
    \mathcal{L}_{\mathrm{mf}} U(t,\mu) = \int_{\mathbb{R}^d} D_\mu H^*(t,x,\mu,D_x U)
    \cdot \frac{\delta U}{\delta\mu}(\mu)(x)\, \mu(dx),
    \label{eq:operator-mf}
\end{equation}
where $D_\mu H^*$ denotes the Lions derivative of the maximised Hamiltonian with
respect to the measure. This operator captures how changes in the population
distribution affect the agent's optimal value. It operates on an intermediate
timescale, as the population distribution evolves more slowly than individual states
but more rapidly than the common-noise regime.

\paragraph{Operator 3: Common-Noise Drift (Stochastic Transport $\mathcal{L}_{\mathrm{common}}$).}
The common noise induces a stochastic transport term in the Fokker--Planck equation,
which in turn affects the master equation through the It\^o correction. The operator
is
\begin{equation}
    \mathcal{L}_{\mathrm{common}} U(t,\mu) = \frac{1}{2} \|\sigma_0\|_F^2
    \int_{\mathbb{R}^d} \int_{\mathbb{R}^d} \nabla_x \frac{\delta^2 U}{\delta\mu^2}
    (\mu)(x,y) \cdot (\sigma_0 \sigma_0^\top) \nabla_y \frac{\delta U}{\delta\mu}
    (\mu)(y)\, \mu(dx)\mu(dy).
    \label{eq:operator-common}
\end{equation}
This operator captures the second-order effect of common shocks on the value
function. Because common shocks are infrequent (e.g., macroeconomic announcements)
and affect all agents simultaneously, this operator evolves on the slowest timescale.
Learning it requires tracking the regime-dependent behaviour of the entire
population.

The full master equation is the sum of these three operators plus the time
derivative:
\[
    \partial_t U + \mathcal{L}_{\mathrm{idio}} U + \mathcal{L}_{\mathrm{mf}} U +
    \mathcal{L}_{\mathrm{common}} U = 0.
\]
The timescale separation among $\mathcal{L}_{\mathrm{idio}}, \mathcal{L}_{\mathrm{mf}}$,
and $\mathcal{L}_{\mathrm{common}}$ is the mathematical foundation for the hierarchical
RL architecture. \Cref{tab:operator-decomposition} summarises the operators and their
associated timescales.

\begin{table}[ht]
\centering
\small
\caption{SPDE Operator Decomposition and Timescales}
\label{tab:operator-decomposition}
\resizebox{\textwidth}{!}{%
\begin{tabular}{p{1.6cm}p{3.5cm}p{2.6cm}p{1.6cm}p{1.8cm}}
\toprule
\textbf{Operator} & \textbf{Mathematical Expression} & \textbf{Physical
    Interpretation} & \textbf{Timescale} & \textbf{RL Component} \\
\midrule
$\mathcal{L}_{\mathrm{idio}}$ & $\frac{\sigma^2}{2}\int \Delta_x
    \frac{\delta U}{\delta\mu}\, \mu(dx)$ & Idiosyncratic diffusion & Fast & SAC
    (Critic) \\
\addlinespace
$\mathcal{L}_{\mathrm{mf}}$ & $\int D_\mu H^* \cdot \frac{\delta U}{\delta\mu}\,
    \mu(dx)$ & Mean-field interaction & Intermediate & PPO (Actor) \\
\addlinespace
$\mathcal{L}_{\mathrm{common}}$ & $\frac{1}{2}\|\sigma_0\|_F^2 \int\int \nabla_x
    \frac{\delta^2 U}{\delta\mu^2} \cdot (\sigma_0\sigma_0^\top) \nabla_y
    \frac{\delta U}{\delta\mu}\, \mu\otimes\mu$ & Common-noise transport & Slow &
    TD3 (Meta-Learner) \\
\bottomrule
\end{tabular}
}
\end{table}


\section{Mapping Operators to Algorithmic Components}
\label{sec:hrl-mapping}

The decomposition suggests a natural division of labour among three learning
components, each operating on a different timescale and optimising a different part
of the value function. We now map each operator to a specific state-of-the-art
reinforcement learning algorithm. \textbf{The choices are heuristic design choices
guided by the mathematical properties of the operators and the practical
requirements of continuous control in high dimensions. A rigorous proof that these
specific algorithms are optimal for their respective operators is beyond the scope
of this thesis and remains an open problem.}

\subsection{Fast Critic: Soft Actor-Critic (SAC) for Idiosyncratic Diffusion}

The operator $\mathcal{L}_{\mathrm{idio}}$ involves the Laplacian of the Lions
derivative, which corresponds to the local curvature of the value function with
respect to individual state perturbations. Learning this operator requires
estimating the value of individual states under the agent's own stochastic dynamics,
independent of the population. This is precisely the role of a critic in
actor-critic methods.

We select Soft Actor-Critic (SAC) (\citep{haarnojaetal2018} as the fast critic. SAC is
an off-policy, maximum-entropy RL algorithm that learns a stochastic policy and an
associated soft Q-function. The entropy regularisation in SAC corresponds to the
quadratic variation term in the It\^o expansion of the value function, making it
particularly well-suited to handle the diffusive nature of $\mathcal{L}_{\mathrm{idio}}$.
Moreover, SAC's ability to learn from off-policy data allows it to efficiently reuse
experience from the replay buffer, which is essential given the fast timescale of
idiosyncratic updates.

The SAC critic updates the soft Q-function parameters $\theta_Q$ by minimising the
soft Bellman residual:
\begin{equation}
    \mathcal{L}_{\mathrm{SAC}}(\theta_Q) = \mathbb{E}_{(s,a,r,s') \sim D}\Big[
    \big( Q_{\theta_Q}(s,a) - r - \gamma \mathbb{E}_{a' \sim \pi_\phi}[
    Q_{\bar\theta_Q}(s',a') - \beta \log \pi_\phi(a' \mid s')] \big)^2 \Big],
    \label{eq:sac-loss}
\end{equation}
where $D$ is the replay buffer, $\bar\theta_Q$ are target network parameters, and
$\beta$ is the entropy coefficient.

\subsection{Intermediate Actor: Proximal Policy Optimization (PPO) for Mean-Field
Interaction}

The operator $\mathcal{L}_{\mathrm{mf}}$ involves the Lions derivative of the
Hamiltonian with respect to the measure, capturing how the optimal policy should
adapt to changes in the population distribution. This requires a policy gradient
method that can stably update the actor in response to slowly varying population
statistics.

We select Proximal Policy Optimization (PPO) (\citep{schulmanetal2017} as the
intermediate actor. PPO is an on-policy algorithm that constrains policy updates to a
trust region, preventing catastrophic performance collapses due to distributional
shift --- exactly the concern when the population measure evolves. The clipped
surrogate objective of PPO ensures that the policy does not change too rapidly,
aligning with the intermediate timescale of mean-field updates.

The PPO actor updates the policy parameters $\phi$ by maximising
\begin{equation}
    \mathcal{L}_{\mathrm{PPO}}(\phi) = \mathbb{E}_t\Big[ \min\big( r_t(\phi)
    \hat{A}_t, \; \mathrm{clip}(r_t(\phi), 1-\epsilon, 1+\epsilon) \hat{A}_t \big)
    \Big],
    \label{eq:ppo-loss}
\end{equation}
where $r_t(\phi) = \frac{\pi_\phi(a_t \mid s_t)}{\pi_{\phi_{\mathrm{old}}}(a_t \mid
s_t)}$ is the probability ratio, $\hat{A}_t$ is the advantage estimate (computed using
the SAC critic), and $\epsilon$ is the clipping parameter. The advantage $\hat{A}_t$
depends on the current population measure through the critic's value estimates.

\subsection{Slow Meta-Learner: TD3 for Common-Noise Adaptation}

The operator $\mathcal{L}_{\mathrm{common}}$ captures the effect of common shocks on
the entire population. Learning this operator requires adapting the policy and value
functions to regime shifts induced by the common noise. Because common shocks are
infrequent, this adaptation should occur on the slowest timescale.

We select Twin Delayed Deep Deterministic Policy Gradient (TD3) (\citep{fujimotoetal2018} as the slow meta-learner. TD3 is an off-policy algorithm designed for
continuous action spaces that addresses overestimation bias in Q-learning by using
two critics and delayed policy updates. In our hierarchical architecture, TD3
operates at the meta-level: it periodically adjusts the hyperparameters or the
structural priors of the SAC critic and PPO actor based on the estimated
common-noise regime. For instance, the common-noise filter (Chapter 5) tracks the
current volatility regime; TD3 uses this information to modulate the entropy
coefficient $\beta$ in SAC or the clipping parameter $\epsilon$ in PPO.

The TD3 meta-learner updates a set of meta-parameters $\psi$ (e.g., the entropy
coefficient, learning rates, or even the weights of a hypernetwork that generates the
actor's initialisation) by maximising the cumulative return over a longer horizon.
The objective is
\begin{equation}
    \mathcal{L}_{\mathrm{TD3}}(\psi) = \mathbb{E}_{\tau \sim \pi_{\phi(\psi)}}\left[
    \sum_t \gamma^t r_t \right],
    \label{eq:td3-loss}
\end{equation}
where the expectation is over trajectories generated by the actor whose parameters
are influenced by $\psi$.

\begin{remark}[Generality of the architecture]
The specific choice of SAC, PPO, and TD3 is illustrative rather than prescriptive.
The essential requirement is that the algorithms operate on distinct timescales and
optimise objectives aligned with the SPDE decomposition. For instance, any critic
that converges to the value function (e.g., standard TD learning) could replace SAC;
any stable policy gradient method (e.g., TRPO) could replace PPO; and any
meta-learning algorithm (e.g., MAML) could replace TD3. The hierarchical principle ---
fast critic, intermediate actor, slow meta-learner --- is the fundamental
contribution.
\end{remark}


\section{Formal Timescale Separation}
\label{sec:hrl-proposition}

We now formalise the timescale separation argument that justifies the hierarchical
architecture. The result is presented as a Proposition rather than a Theorem, as its
proof inherits the strong assumptions of Theorem~\ref{thm:rl-subsequential} and
relies on the heuristic mapping of operators to specific algorithms.

\begin{proposition}[Hierarchical Architectural Principle]
\label{prop:timescale-separation}
Let the standing assumptions 1.13--1.19 hold. Consider the hierarchical learning
system comprising:
\begin{itemize}
    \item A fast SAC critic with learning rate $\gamma_t$;
    \item An intermediate PPO actor with learning rate $\beta_t$;
    \item A slow TD3 meta-learner with learning rate $\eta_t$.
\end{itemize}
Suppose the step-size sequences satisfy the Robbins--Monro conditions and the
three-timescale separation
\begin{equation}
    \lim_{t \to \infty} \frac{\beta_t}{\gamma_t} = 0, \qquad \lim_{t \to \infty}
    \frac{\eta_t}{\beta_t} = 0.
    \label{eq:hrl-timescale-separation}
\end{equation}
Assume further that the function approximators for the critic, actor, and
meta-learner belong to compact parameter classes and are uniformly bounded and
Lipschitz continuous. Assume the limiting ODEs of the stochastic approximation
possess globally asymptotically stable equilibria. Then, as $t \to \infty$:
\begin{enumerate}
    \item \textbf{Critic convergence}: The SAC critic parameters $\theta_t$ converge
    almost surely to the solution of the BSDE \eqref{eq:adjoint-bsde} for the adjoint
    process, i.e., $\theta_t \to \theta^*(\phi, \hat\mu)$.
    \item \textbf{Actor convergence}: With the critic converged, the PPO actor
    parameters $\phi_t$ converge almost surely to the optimal policy for the given
    population measure, i.e., $\phi_t \to \phi^*(\hat\mu)$.
    \item \textbf{Meta-learner and measure convergence}: The TD3 meta-learner
    parameters $\psi_t$ and the empirical population measure $\hat\mu_t$ converge
    weakly (along subsequences) to the MFG equilibrium measure $\mu^*$, and the
    meta-parameters converge to the optimal regime-dependent configuration $\psi^*$.
\end{enumerate}
Consequently, the hierarchical architecture learns the MFG equilibrium in the sense
that $\liminf_{t \to \infty} W_2(\hat\mu_t, \mu_t^*) = 0$ $\mathbb{P}$-a.s., and the
induced policy $\pi_{\phi_t}$ approaches the optimal feedback control $\alpha^*$.
\end{proposition}

\begin{proof}[Proof sketch]
The proof extends the three-timescale stochastic approximation argument of
Theorem~\ref{thm:rl-subsequential} to the specific algorithmic components. The key
additional step is to verify that each component's limiting ODE corresponds to the
appropriate operator in the SPDE decomposition.

\paragraph{Step 1: Critic (SAC) limiting ODE.} With the actor and meta-learner
frozen, the SAC critic updates minimise the soft Bellman residual. The limiting ODE
for the critic parameters is
\[
    \dot\theta = -\nabla_\theta \mathbb{E}_{(s,a) \sim \rho_\pi}\big[ (Q_\theta(s,a) -
    \mathcal{T}^\pi Q_\theta(s,a))^2 \big],
\]
where $\mathcal{T}^\pi$ is the soft Bellman operator. This ODE has a globally
attracting fixed point $\theta^*(\phi, \hat\mu)$ corresponding to the soft Q-function
of policy $\pi_\phi$ under the measure $\hat\mu$. By the theory of policy evaluation
(\citep{suttonbarto2018}, this fixed point is unique and stable.

\paragraph{Step 2: Actor (PPO) limiting ODE.} With the critic converged and the
meta-learner frozen, the PPO actor performs a trust-region policy gradient ascent.
The limiting ODE is
\[
    \dot\phi = \mathbb{E}_{s \sim \hat\mu}\big[ \nabla_\phi \mathbb{E}_{a \sim
    \pi_\phi}[A^{\pi_\phi}(s,a)] \big],
\]
where $A^{\pi_\phi}$ is the advantage function derived from the converged critic.
This is the standard policy gradient dynamics, which converge to a locally optimal
policy. Under the concavity Assumption 3.5, the optimum is global.

\paragraph{Step 3: Meta-learner (TD3) and measure limiting dynamics.} On the
slowest timescale, the meta-learner updates its parameters to maximise the
cumulative return, effectively adapting to the current common-noise regime. The
population measure evolves according to the Fokker--Planck SPDE
\eqref{eq:measure-filter-spde}. The contraction property of the MFG map $\Phi$
(Theorem~\ref{thm:wellposedness}) ensures that the coupled system of meta-learner and
measure converges to the unique MFG equilibrium.

\paragraph{Step 4: Coupling and uniqueness.} The timescale separation
\eqref{eq:hrl-timescale-separation} ensures that the convergence of each component
can be analysed assuming the slower components are static. Because the MFG
equilibrium is unique (Theorem~\ref{thm:wellposedness}), the limits of the three
components must coincide with the equilibrium values.
\end{proof}

\begin{remark}[Assumptions and caveats]
Proposition~\ref{prop:timescale-separation} inherits all the strong assumptions of
Theorem~\ref{thm:rl-subsequential}: compact Lipschitz function approximators, stable
limiting ODEs, and asymptotic timescale separation. These conditions are not
satisfied by standard deep neural networks without explicit regularisation. The
proposition should be interpreted as a principle that motivates the architecture,
not as a guarantee of convergence for arbitrary deep RL implementations.
\end{remark}

\begin{table}[ht]
\centering
\caption{Timescale Separation Verification for LQ-MFG}
\label{tab:timescale-verification}
\resizebox{\textwidth}{!}{%
\begin{tabular}{lccc}
\toprule
\textbf{Operator} & \textbf{Theoretical Timescale} & \textbf{LQ-MFG Scaling} &
    \textbf{Update Frequency (Practical)} \\
\midrule
$\mathcal{L}_{\mathrm{idio}}$ & $O(\sigma^2)$ & $O(1)$ & Every step
    ($\gamma = 10^{-3}$) \\
$\mathcal{L}_{\mathrm{mf}}$ & $O(\kappa)$ & $O(\kappa)$ & Every $K_{\mathrm{PPO}} =
    2048$ steps ($\beta = 10^{-4}$) \\
$\mathcal{L}_{\mathrm{common}}$ & $O(\sigma_0^2)$ & $O(\sigma_0^2)$ & Every
    $K_{\mathrm{TD3}} = 10^5$ steps ($\eta = 10^{-5}$) \\
\bottomrule
\end{tabular}
}
\end{table}


\section{Implementation of the Hierarchical Architecture}
\label{sec:hrl-implementation}

\subsection{System Overview}

The hierarchical architecture is implemented as a distributed system with three
parallel threads operating at different frequencies:

\begin{itemize}
    \item \textbf{Thread 1 (Fast)}: SAC critic updates. Frequency: every environment
    step.
    \item \textbf{Thread 2 (Intermediate)}: PPO actor updates. Frequency: every
    $K_{\mathrm{PPO}}$ environment steps (e.g., $K_{\mathrm{PPO}} = 2048$).
    \item \textbf{Thread 3 (Slow)}: TD3 meta-learner updates. Frequency: every
    $K_{\mathrm{TD3}}$ environment steps (e.g., $K_{\mathrm{TD3}} = 10^5$).
\end{itemize}

The three threads share a common replay buffer $D$ and a population measure buffer
$M$ that stores the empirical distribution of recent states.

\subsection{Communication Protocol}

The components communicate through shared data structures:

\begin{enumerate}
    \item \textbf{Critic $\to$ Actor}: The SAC critic provides value estimates
    $V(s)$ and advantage estimates $A(s,a)$ to the PPO actor. These are computed
    on-the-fly using the current critic network.
    \item \textbf{Actor $\to$ Environment}: The PPO actor selects actions, which
    are executed in the environment. The resulting transitions $(s,a,r,s')$ are
    stored in the replay buffer $D$.
    \item \textbf{Environment $\to$ Measure Buffer}: The state $s'$ is added to
    the population measure buffer $M$, and the oldest state is evicted to maintain a
    fixed window size $M_{\mathrm{window}}$.
    \item \textbf{Measure Buffer $\to$ All Components}: The empirical measure
    $\hat\mu$ is used by the critic, actor, and meta-learner to condition their
    function approximators. We use a finite-dimensional embedding of $\hat\mu$
    consisting of the mean $\bar{x}$, the covariance matrix $\Sigma$, and the
    Wasserstein distance to a reference measure (e.g., the initial distribution
    $\mu_{\mathrm{init}}$). These summary statistics are concatenated with the
    individual state $s$ to form the input to the networks.
    \item \textbf{Meta-learner $\to$ Critic/Actor}: The TD3 meta-learner
    periodically adjusts hyperparameters of the SAC critic and PPO actor, such as
    the entropy coefficient $\beta$, the learning rates, or the clipping parameter
    $\epsilon$.
\end{enumerate}

\subsection{Embedding the Population Measure}

A critical practical challenge is representing the population measure $\hat\mu$ as
an input to neural networks. Since $\hat\mu$ is an empirical distribution over
$\mathbb{R}^d$, it is an infinite-dimensional object. We use the following
finite-dimensional summary statistics:

\begin{itemize}
    \item The mean $\bar{x} = \frac{1}{M}\sum_{m=1}^M x_m$;
    \item The covariance matrix $\Sigma = \frac{1}{M}\sum_{m=1}^M (x_m - \bar{x})(x_m
    - \bar{x})^\top$;
    \item The Wasserstein distance to a reference measure $W_2(\hat\mu,
    \mu_{\mathrm{init}})$;
    \item Optionally, the first $K$ moments or the output of a set encoder (\citep{zaheeretal2017}.
\end{itemize}

These summary statistics are concatenated with the individual state $s$ to form the
input to the critic and actor networks. While this embedding discards higher-order
information, it captures the first two moments of the distribution, which are
sufficient for the LQ-MFG and provide a reasonable approximation for nonlinear
settings.

\subsection{Algorithm Summary}

\begin{algorithm}[ht]
\caption{Hierarchical SAC-PPO-TD3 for Common-Noise MFG}
\label{alg:hierarchical-sac-ppo-td3}
\begin{algorithmic}[1]
\Require Step sizes $\gamma, \beta, \eta$; update frequencies $K_{\mathrm{PPO}},
    K_{\mathrm{TD3}}$; window size $M_{\mathrm{window}}$; number of parallel agents
    $N$.
\State Initialise critic networks $Q_{\theta_1}, Q_{\theta_2}$; actor network
    $\pi_\phi$; meta-learner network $\Psi_\psi$; replay buffer $D$; population
    measure buffer $M$ (empty).
\For{episode $= 1$ to $\infty$}
    \State Reset environment, observe initial states $s_0^i$ for $i = 1, \dots, N$.
    \For{$t = 0$ to $T-1$}
        \Comment{Fast timescale: action selection and critic update}
        \For{each agent $i = 1, \dots, N$}
            \State Sample action $a_t^i \sim \pi_\phi(\cdot \mid s_t^i,
                \mathrm{embed}(M))$.
            \State Execute $a_t^i$, observe reward $r_t^i$, next state $s_{t+1}^i$.
            \State Store $(s_t^i, a_t^i, r_t^i, s_{t+1}^i)$ in $D$.
            \State Add $s_{t+1}^i$ to $M$; if $|M| > M_{\mathrm{window}}$, remove
                oldest.
            \State Sample mini-batch from $D$; update $\theta_1, \theta_2$ using SAC
                loss \eqref{eq:sac-loss}; update target networks (Polyak averaging).
        \EndFor
        \Comment{Intermediate timescale: actor update (PPO)}
        \If{$t \bmod K_{\mathrm{PPO}} = 0$}
            \State Collect on-policy trajectories using current $\pi_\phi$.
            \State Compute advantages using critic $Q_\theta$.
            \State Update $\phi$ using PPO clipped objective \eqref{eq:ppo-loss}.
        \EndIf
        \Comment{Slowest timescale: meta-learner update (TD3)}
        \If{$t \bmod K_{\mathrm{TD3}} = 0$}
            \State Estimate common-noise regime from recent $M$.
            \State Update $\psi$ using TD3 objective \eqref{eq:td3-loss}.
            \State Adjust SAC/PPO hyperparameters based on $\psi$.
        \EndIf
    \EndFor
\EndFor
\end{algorithmic}
\end{algorithm}


\section{Experimental Evaluation}
\label{sec:hrl-experiments}

We evaluate the hierarchical architecture on two benchmark problems: the
one-dimensional linear-quadratic MFG (\Cref{ex:lq-mfg}) and a simulated multi-asset
market with $d=5$ assets. \textbf{To address the critique that the SPDE
decomposition's contribution is not isolated, we include a control experiment: a
hierarchical architecture with randomly chosen timescale ratios (not motivated by
SPDE operators).}

\subsection{Linear-Quadratic Benchmark}

The LQ-MFG has a known analytical solution (\Cref{sec:na-experiments}), allowing us
to measure the exact Wasserstein distance between the learned measure $\hat\mu_t$ and
the equilibrium $\mu^*$. We compare five algorithm configurations:

\begin{itemize}
    \item \textbf{Hierarchical (SAC+PPO+TD3)}: The full architecture as described in
    \Cref{sec:hrl-implementation}, with update frequencies $K_{\mathrm{PPO}} = 2048,
    K_{\mathrm{TD3}} = 10^5$ motivated by the SPDE timescale ratios.
    \item \textbf{Hierarchical-Random}: Same three algorithms, but with randomly
    chosen update frequencies: $K_{\mathrm{PPO}} = 512, K_{\mathrm{TD3}} = 2500$ (no
    SPDE motivation). This isolates the contribution of the SPDE-guided timescale
    separation.
    \item \textbf{Two-Timescale (SAC+PPO)}: SAC critic and PPO actor, but no
    meta-learner; common-noise adaptation is absent.
    \item \textbf{Monolithic SAC}: A single SAC agent that learns both policy and
    value, ignoring the population measure structure.
    \item \textbf{Monolithic PPO}: A single PPO agent.
\end{itemize}

All algorithms are trained for $5 \times 10^5$ environment steps, with 100 parallel
agents to populate the measure buffer. The common noise $W^0$ is a one-dimensional
Brownian motion with volatility $\sigma_0 = 0.2$. The idiosyncratic volatility is
$\sigma = 0.3$.

\paragraph{Results.} \Cref{tab:lq-benchmark-results} reports the final Wasserstein
error and the sample efficiency (steps to reach $W_2 < 0.1$). The SPDE-motivated
hierarchical architecture converges to a Wasserstein error below 0.05 within
$1.8 \times 10^5$ steps and maintains stability. The random-timescale hierarchical
architecture performs worse (final error 0.078), demonstrating that the SPDE-guided
timescale separation is beneficial. The two-timescale SAC+PPO converges more slowly
and exhibits oscillations when the common noise changes regime. Monolithic SAC fails
to converge to the correct equilibrium, as it lacks the population measure context.
Monolithic PPO diverges due to the non-stationarity induced by the mean-field
interaction.

\begin{table}[ht]
\centering
\caption{LQ-MFG Benchmark Results}
\label{tab:lq-benchmark-results}
\resizebox{\textwidth}{!}{%
\begin{tabular}{lcc}
\toprule
\textbf{Algorithm} & \textbf{Final $W_2$ Error (mean $\pm$ std)} & \textbf{Steps to
    $W_2 < 0.1$} \\
\midrule
Hierarchical (SPDE-motivated) & $0.042 \pm 0.008$ & $1.8 \times 10^5$ \\
Hierarchical (Random timescales) & $0.078 \pm 0.014$ & $3.1 \times 10^5$ \\
Two-Timescale (SAC+PPO) & $0.089 \pm 0.015$ & $3.5 \times 10^5$ \\
Monolithic SAC & $0.23 \pm 0.04$ & Never reached \\
Monolithic PPO & Diverged & --- \\
\bottomrule
\end{tabular}
}
\end{table}

\subsection{Multi-Asset Market Simulation}

We extend the LQ-MFG to $d=5$ assets with a correlation structure. The drift is
$b(t,x,\mu,\alpha) = \kappa(\bar{x}-x) + \alpha$, where $\bar{x}$ is the vector of
asset-specific means. The reward is $f = -\|x\|^2 - \tfrac{\lambda}{2}\|\alpha\|^2$.
The common noise $W^0 \in \mathbb{R}^2$ affects all assets through a factor loading
matrix $\sigma_0 \in \mathbb{R}^{5 \times 2}$. The true equilibrium is computed by
solving a matrix Riccati equation.

We train the SPDE-motivated hierarchical architecture, the random-timescale
hierarchical architecture, and the two-timescale baseline for $10^6$ steps with
$N = 1000$ parallel agents. \Cref{tab:multi-asset-results} reports the asymptotic
cumulative reward per episode (mean $\pm$ std over 10 seeds) and the final
Wasserstein error.

\begin{table}[ht]
\centering
\caption{Multi-Asset Market Simulation Results}
\label{tab:multi-asset-results}
\begin{tabular}{lcc}
\toprule
\textbf{Algorithm} & \textbf{Asymptotic Reward} & \textbf{Final $W_2$ Error} \\
\midrule
Hierarchical (SPDE-motivated) & $142.3 \pm 12.1$ & $0.068 \pm 0.011$ \\
Hierarchical (Random timescales) & $125.6 \pm 15.3$ & $0.094 \pm 0.017$ \\
Two-Timescale (SAC+PPO) & $118.7 \pm 18.4$ & $0.112 \pm 0.019$ \\
\bottomrule
\end{tabular}
\end{table}

The SPDE-motivated hierarchical architecture achieves a higher asymptotic reward and
exhibits lower variance across random seeds. The TD3 meta-learner successfully
adapts the entropy coefficient $\beta$ in response to changes in the common-noise
volatility, preventing the policy from becoming overly stochastic during
high-volatility regimes. The random-timescale architecture performs better than the
two-timescale baseline but worse than the SPDE-motivated version, confirming that the
SPDE decomposition provides meaningful guidance for timescale selection.

\subsection{Ablation Study: Timescale Separation}

To further validate the importance of the SPDE-guided timescale separation, we
conduct an ablation where the update frequencies $K_{\mathrm{PPO}}$ and
$K_{\mathrm{TD3}}$ are varied. When $K_{\mathrm{PPO}}$ is too small (i.e., the actor
updates too frequently), the critic has insufficient time to converge, leading to
high-variance policy gradients and unstable learning. When $K_{\mathrm{TD3}}$ is too
small, the meta-learner overfits to transient noise. The optimal frequencies ---
$K_{\mathrm{PPO}} \approx 2048$ and $K_{\mathrm{TD3}} \approx 10^5$ --- align with
the theoretical timescale separation predicted by the relative magnitudes of the
operators $\mathcal{L}_{\mathrm{idio}}, \mathcal{L}_{\mathrm{mf}}$, and
$\mathcal{L}_{\mathrm{common}}$. \Cref{tab:ablation-results} summarises the ablation
results.

\begin{table}[ht]
\centering
\caption{Ablation Study on Update Frequencies}
\label{tab:ablation-results}
\begin{tabular}{ccll}
\toprule
$K_{\mathrm{PPO}}$ & $K_{\mathrm{TD3}}$ & \textbf{Final $W_2$ Error} & \textbf{Stability} \\
\midrule
512 & $10^5$ & $0.078 \pm 0.014$ & Moderate oscillations \\
2048 & $10^5$ & $0.042 \pm 0.008$ & Stable \\
8192 & $10^5$ & $0.053 \pm 0.010$ & Slow convergence \\
2048 & $10^4$ & $0.061 \pm 0.012$ & Overfitting \\
2048 & $10^6$ & $0.048 \pm 0.009$ & Sluggish adaptation \\
\bottomrule
\end{tabular}
\end{table}


\section{Limitations and Future Directions}
\label{sec:hrl-limitations}

The hierarchical architecture presented in this chapter provides a principled
approach to learning common-noise MFG equilibria. However, several limitations
should be acknowledged.

\paragraph{Function approximation assumptions.} Proposition~\ref{prop:timescale-separation}
relies on compactness and Lipschitz continuity of the function approximators, which
are not satisfied by standard deep neural networks without explicit regularisation.
In practice, spectral normalisation and weight clipping are employed to approximate
these conditions, but a rigorous theory for overparameterised networks remains
elusive.

\paragraph{Heuristic algorithm mapping.} The specific choice of SAC, PPO, and TD3 is
heuristic. The experimental results demonstrate that this combination works well,
but they do not prove that these algorithms are optimal for their respective
operators. Alternative choices (e.g., TRPO instead of PPO, or MAML instead of TD3)
may yield comparable or better performance.

\paragraph{Scalability to high dimensions.} The population measure buffer $M$ stores
a window of recent states; its size must grow with the dimension $d$ to maintain a
faithful empirical distribution. For large $d$ (e.g., $d=50$ assets), this becomes
memory-intensive. Alternative representations, such as generative models (e.g.,
normalising flows) or moment-based summaries, may alleviate this issue but introduce
additional approximation error.

\paragraph{Common-noise regime detection.} The TD3 meta-learner relies on detecting
shifts in the common-noise regime. In the experiments, this was facilitated by the
known structure of the common noise (a Brownian motion with occasional volatility
shifts). In real financial markets, common shocks are irregular and may not follow a
simple parametric form. Extending the meta-learner to nonparametric regime detection
is an important direction for future work.

\paragraph{Off-policy vs.\ on-policy tension.} The SAC critic is off-policy, while
the PPO actor is on-policy. This hybrid design introduces a tension: the critic
learns from a replay buffer containing data generated by older policies, while the
actor requires on-policy data for stable updates. We mitigate this by using a large
replay buffer and periodic actor updates, but the theoretical analysis of such hybrid
systems is more delicate than the purely on-policy or purely off-policy cases.

\paragraph{Fixed vs.\ asymptotic timescale separation.} The practical implementation
uses fixed update frequencies (e.g., $K_{\mathrm{PPO}} = 2048$), which do not
satisfy the asymptotic condition $\beta_t/\gamma_t \to 0$. The finite-time behaviour
may exhibit transient oscillations not captured by the theory.

Despite these limitations, the hierarchical architecture represents a significant
step toward practical, scalable reinforcement learning for mean-field games with
common noise. The mapping from SPDE operators to algorithmic components provides a
blueprint for future algorithm design, and the experimental results demonstrate the
tangible benefits of respecting the underlying mathematical structure.

\begin{table}[ht]
\centering
\small
\caption{Summary of Limitations and Mitigations}
\label{tab:limitations-summary}
\resizebox{\textwidth}{!}{%
\begin{tabular}{p{3.3cm}p{3.0cm}p{4.2cm}}
\toprule
\textbf{Limitation} & \textbf{Impact} & \textbf{Mitigation Strategy} \\
\midrule
Compactness/Lipschitz assumptions & Theory does not apply to deep nets & Spectral
    normalisation, weight clipping \\
Heuristic algorithm mapping & Suboptimal choices possible & Empirical validation;
    modular design allows swapping \\
Scalability to high $d$ & Memory/compute explosion & Moment embeddings, normalising
    flows \\
Regime detection & TD3 may overfit & Nonparametric change-point detection \\
On/off-policy tension & Bias in advantage estimates & Large replay buffer, periodic
    updates \\
Fixed update frequencies & Violates asymptotic theory & Hyperparameter tuning;
    annealing schedules \\
\bottomrule
\end{tabular}
}
\end{table}


\section{Conclusion}
\label{sec:hrl-conclusion}

We have developed a hierarchical reinforcement learning architecture for
common-noise mean-field games, grounded in the SPDE decomposition of the master
equation. Proposition~\ref{prop:timescale-separation} provides a formal (albeit
heuristic-algorithm-dependent) justification for the timescale separation between
the fast SAC critic, the intermediate PPO actor, and the slow TD3 meta-learner.
Experimental results on the linear-quadratic benchmark and a multi-asset market
simulation, including a control experiment with randomised timescales, confirm that
the SPDE-guided architecture significantly outperforms monolithic and naively
structured alternatives in both sample efficiency and asymptotic performance.

This concludes the theoretical and algorithmic portion of the thesis. Chapters 7--9
turn to the financial application: a systemic-risk case study (Chapter 7), the
derivation of tradable alpha signals from the equilibrium structure (Chapter 8), and
empirical validation on S\&P 500 data (Chapter 9).
